\documentclass{article}
\usepackage{graphicx}
\usepackage[russian]{babel}
\usepackage[14pt]{extsizes}
\usepackage[utf8x]{inputenc}
\linespread{1.5}
\usepackage{indentfirst}
\setlength{\parindent}{12.5mm}
\usepackage{multirow}
\usepackage{tabularx}
\usepackage{pdflscape}
\usepackage{geometry}
\geometry{top=2cm}
\geometry{right=1cm}
\geometry{left=3cm}
\geometry{bottom=2cm}
\usepackage{amsthm}
\usepackage{amsmath}
\usepackage{amssymb}
\usepackage{amsfonts}
\usepackage{graphicx}
\usepackage{xcolor}
\usepackage{xlop}
\usepackage{tocloft}
\usepackage{algorithm}
\usepackage{algpseudocode}
\usepackage{float} 
\usepackage{supertabular}
\usepackage{booktabs}
\usepackage{multicol}
\usepackage{caption}

\begin{document}

\section{Теоретический минимум}

Формат модуля: какая-то теория, теоремы, определения, умные слова, все что угодно. Каждую мало-мальски тему выделяем в отдельную подсекцию.

\subsection{Какие-то понятия из предметной области}
\textbf{Опцион} --  договор, по которому потенциальный покупатель или потенциальный продавец базового актива (товара, ценной бумаги) получает право, но не обязательство совершить покупку или продажу данного актива по заранее оговорённой цене. Рассматриваются европейские опционы, т.е. опционы с фиксированной датой исполнения (экспирации).

Опцион \textbf{колл} -- опцион на покупку. Предоставляет покупателю опциона право купить базовый актив по фиксированной цене.

Опцион \textbf{пут} -- опцион на продажу. Предоставляет покупателю опциона право продать базовый актив по фиксированной цене.
\[
\textbf{ТУТ ЯВНО НУЖНО ЧТО-ТО  ЕЩЕ}
\]
\newpage
\subsection{Модель Блека-Шоулза}
\textbf{Цена опциона колл}:
\[
C = SN(d_1) - Ke^{-rT}N(d_2),
\]

\textbf{Цена опциона пут}:
\[
P= Ke^{-rT}N(d_2) - SN(-d_1)
\]
\[
d_1 = \frac{ln\frac{S}{K}+ \big(r + \frac{\sigma^2}{2}\big)T}{\sigma\sqrt{T}}, \quad d_2 = d_1 - \sigma\sqrt{T}
\]
\[
\textbf{МБ МОЖНО КАК-ТО БОЛЕЕ КРАСИВО ЗАПИСАТЬ ФОРМУЛЫ}
\]
Обозначения:
\begin{itemize}
    \item $C$ -- цена опциона call;
    \item $S$ -- текущая цена базисного актива (спот);
    \item $N(x)$ -- кумулятивная функция распределения стандартного нормального распределения;
    \item $K$ -- цена исполнения опциона (страйк);
    \item $r$ -- безрисковая процентная ставка;
    \item $T$ -- время до экспирации опциона (считается в годах);
    \item $\sigma$ -- волатильность доходности базисного актива.
\end{itemize}
\[
\textbf{ТУТ ЯВНО НУЖНО ЧТО-ТО  ЕЩЕ}
\]
\newpage
\subsection{А что вообще мы делаем?}

\newpage
\section{Какие-то более технические выкладки}

Как мне кажется, тут идеологически должно быть что-то типо двух частей: первая часть о всех численных методах, которые используются с каким-то разумным количеством теории и каким-то выводом к чему приходим в итоге. Можно сразу бахнуть реализацию прямо в этих параграфах, можно в целом не бахать. Вторая часть про передачу данных и эзернетус. 
\subsection{Как посчитать $\Phi(x)$?}
\textbf{Нормальное распределение}, также называемое распределением Гаусса или Гаусса-Лапласа, или колоколообразная кривая -- непрерывное распределение вероятностей с пиком в центре и симметричными боковыми сторонами, которое в одномерном случае задаётся функцией плотности вероятности, совпадающей с функцией Гаусса:
\[
f(x) = \frac{1}{\sqrt{2\pi}\sigma}\operatorname{exp}\bigg({-\frac{1}{2}\Big(\frac{x-\mu}{\sigma}\Big)^2}\bigg)
\]

\textbf{Стандартным нормальным распределением} называется нормальное распределение с математическим ожиданием $\mu = 0$ и стандартным отклонением $\sigma = 1$. Его функция плотности веротяности будет:
\[
\varphi(x) = \frac{1}{\sqrt{2\pi}}\operatorname{exp}\bigg({-\frac{x^2}{2}}\bigg)
\]

\textbf{Функция распределения стандартного нормального распределения} (нормальное интегральное распределение) обычно обозначается заглавной греческой буквой $\Phi$, в русскоязычной литературе называется функцией Лапласа и представляет собой интеграл:
\[
\Phi(x) = \frac{1}{\sqrt{2\pi}}\int\limits_{-\infty}^{x}e^{-\frac{t^2}{2}}dt
\]

Итого:
\[
N(d_1) = \frac{1}{\sqrt{2\pi}}\int\limits_{-\infty}^{d_1}e^{-\frac{t^2}{2}}dt
\]
\[
N(d_2) = \frac{1}{\sqrt{2\pi}}\int\limits_{-\infty}^{d_2}e^{-\frac{t^2}{2}}dt
\]

Нужно как-то уметь вычислять это интеграл, так как он вообще говоря не берущийся...
\[
\textbf{ТУТ ЯВНО НУЖНО ЧТО-ТО  ЕЩЕ}
\]
\subsection{Мультисекция}
\[
\textbf{ТУТ ЯВНО НУЖНО ЧТО-ТО  ЕЩЕ}
\]
\section{Какая-то часть в которой все собирается воедино}
\[
\textbf{ТУТ ЯВНО НУЖНО ЧТО-ТО  ЕЩЕ}
\]
\end{document}
